\subsection{The non-finite-flat branch and a uniform modular budget}
\label{fm7-section}

The condition $p\nmid V$ in Theorem~\ref{fm-residual-budget} can be
removed by retaining a second weight. This is a necessary modular
condition and a dimension estimate; candidate elimination remains an
additional problem.

\begin{theorem}[Both residual-weight branches]\label{fm7-weight-level}
Let $a,b$ be coprime positive integers, let $p>7$ be prime, and suppose
\[
 F(a,b)=VQ^p,\qquad V,Q\in\mathbb Z_{>0},
\]
where $V$ is $p$-power-free. Put $r=v_p(V)$, $R=V/p^r$, and
\begin{equation}\label{fm7-weight}
 k(V,p)=\begin{cases}2,&r=0,\\p+1,&r>0.\end{cases}
\end{equation}
The descended residual representation constructed in
Theorem~\ref{fm-five-levels} arises from a normalized characteristic-zero
newform of weight $k(V,p)$, nebentypus $\psi_3$ of conductor $12$, and
exact prime-to-$p$ level
\begin{equation}\label{fm7-level}
 N=2^e3^2\operatorname{rad}(R),\qquad 2\le e\le6.
\end{equation}
If $r>0$, then $p\equiv1\pmod3$. In both branches, for
$\lambda=\max\{1,\log V\}/p$,
\begin{equation}\label{fm7-level-height}
 \frac{\log N}{p}\le\lambda+\frac{\log576}{p}.
\end{equation}
\end{theorem}

\begin{proof}
The actual Frey curve is the non-CM Q-curve of minimum degree $2$
already studied above. Its multiplicative prime greater than $3$, the
fixed order-four splitting character $\chi$, the odd descent with
determinant $\psi_3$ times the cyclotomic character, and the general
large-image theorem giving residual absolute irreducibility for $p>7$
do not require $p\nmid V$.

Suppose $r>0$. Since $p\mid F(a,b)$, the local calculation of the Frey
curve shows that $p$ splits in $K=\mathbb Q(\sqrt{-3})$, and either
completion is $\mathbb Q_p$. The two minimal discriminant valuations
are $m,2m$, where
\begin{equation}\label{fm7-tate-depth}
 m=r+p\,v_p(Q),\qquad1\le r\le p-1.
\end{equation}
The corresponding $c_4$ valuations are zero, so the $j$ valuations are
$-m,-2m$, neither divisible by $p$. The character $\chi$ is unramified
at $p$, so it leaves the inertia representation and Serre weight
unchanged. Serre's semistable weight formula
\cite[\S2.9, Proposition~5]{Serre1987Modular} therefore gives exactly
$k=p+1$. That formula includes nonsplit multiplicative reduction by
an unramified quadratic extension. If $r=0$, the good-reduction or
finite-flat Tate argument of Theorem~\ref{fm-residual-budget} instead
gives weight $2$, including an inert prime $p$ in the good-reduction
case.

For $q>3$, $q\ne p$, the Tate monodromy modulo $p$ is nonzero precisely
when $p\nmid v_q(F)$. Since $V$ is $p$-power-free, this is equivalent
to $q\mid V$. Thus the prime-to-$p$ conductor factor away from $6$ is
$\operatorname{rad}(R)$. The conductor arguments at $2$ and $3$ are
unchanged: their exponents are $2\le e\le6$ and exactly $2$. Reduction
modulo $p>7$ preserves these finite inertia conductors. This proves
\eqref{fm7-level} for the residual conductor.

The completed strong form of Serre modularity, due to
Khare--Wintenberger with the Kisin input described in their work
\cite{KhareWintenberger2009Serre}, supplies a residual eigensystem at
this Serre weight and conductor. Serre's 1987 paper supplies the local
weight formula and characteristic-zero lifting, not a proof of the
then-conjectural full modularity theorem. Both weights satisfy
$k-1\equiv1\pmod{p-1}$, so the determinant identifies the residual
nebentypus as $\psi_3$. Its Teichm\"uller lift is the same quadratic
character. The characteristic-zero eigensystem lift preserves weight,
level and this character \cite[\S3.1.6]{Serre1987Modular}.
It can be chosen new at exact level $N$: an oldform from a proper divisor
of $N$ would give residual conductor dividing that smaller level,
contradicting the exact conductor just computed. Finally $N\le576V$
gives \eqref{fm7-level-height}.
\end{proof}

\begin{corollary}[Uniform dimension budget]\label{fm7-budget}
For each $p>7$ and each positive $p$-power-free $V$ coprime to $6$, the
sum of the dimensions of the five newspaces in
\eqref{fm7-weight}--\eqref{fm7-level} is at most
\begin{equation}\label{fm7-fixed-budget}
 485(p+1)V^2.
\end{equation}
For every $L\ge0$, their dimension sum over all $1\le V\le e^{Lp}$
satisfying these hypotheses is at most
\begin{equation}\label{fm7-moving-budget}
 485(p+1)e^{3Lp}.
\end{equation}
Each $V$ contributes its one prescribed weight $k(V,p)$. This counts
Galois orbits with their coefficient-field degrees, and includes
residual integers for which no actual seed exists.
\end{corollary}

\begin{proof}
Put $D=\operatorname{rad}(R)$. Since $\gcd(V,6)=1$, the modular-group
index satisfies
\[
 I=[\mathrm{SL}_2(\mathbb Z):\Gamma_0(N)]
   =2N\prod_{q\mid D}(1+q^{-1})\le2ND\le1152V^2.
\]
The factors at $2$ and $3$ have product $2$, and
$1+q^{-1}\le q$ bounds the other factors. The Sturm coefficient bound
\cite{PARIModularForms} bounds the whole cusp space, hence its newspace,
by $1+\lfloor kI/12\rfloor\le1+96(p+1)V^2$. Summing over five
exponents $e$ gives \eqref{fm7-fixed-budget}. For
$H=\lfloor e^{Lp}\rfloor$, use $\sum_{V=1}^H V^2\le H^3$ to obtain
\eqref{fm7-moving-budget}. In particular,
\[
 \frac{\log(485(p+1)e^{3Lp})}{p}
   =3L+\frac{\log485+\log(p+1)}p\longrightarrow0
\]
when $L\to0$ and $p\to\infty$.
\end{proof}

\paragraph{The remaining condition.}
Every actual power profile in this budget gives a congruence, at a place
above $p$, between its descended Frey representation and a member of
the indicated finite newform collection. Its seed also satisfies the
exact inverse-square and positivity conditions above. The missing input
is an exclusion of these congruences for actual integer seeds, uniformly
in moving residual integers and primes. A dimension bound does not
bound seed height or exclude a single congruence.
The nontrivial compression target has $Q>1$. The necessary theorem
also includes $Q=1$; fixed seeds can have such trivial representations
for arbitrarily large $p$, so no blanket exclusion of every profile
is intended.

For example, a theorem forcing every such congruence into a finite
explicitly classified seed list, together with verification that the
list contains no required profiles, would close that branch. No such
theorem is asserted here. A Frey--Mazur type claim turning a congruence
with the boundary curve into an isogeny is likewise an additional
conditional input; finite matching traces do not prove it. The local
boundary shadows are not asserted to be global power profiles.

For $V=1$, exclusion at all sufficiently large prime exponents, combined
with the earlier fixed-exponent finiteness, would settle the pure
second-norm family with exponent at least $3$. It would not settle
arbitrary residual profiles or ABC. Moreover, if $p\mid g$, first
write $V=V_pC^p$ with $V_p$ $p$-power-free. The new profile is
$F=V_p(CQ^{g/p})^p$, and its parameter satisfies
$\lambda_p=\max\{1,\log V_p\}/p\le\max\{1,\log V\}/p$.
It is at most $1/\theta$ times $\max\{1,\log V\}/g$ when
$p\ge\theta g$ for a fixed $\theta>0$; without this condition its
smallness does not follow from the original parameter's smallness.
The independent Kummer point-height route and these remaining modular
branches stay open.
